\documentclass{article}
\usepackage[UTF8]{ctex}
\usepackage{amsmath,amssymb,amsfonts}
\usepackage{enumitem}
\usepackage{geometry}
\geometry{a4paper, margin=1in}
% === 错题本排版优化（不参与题目标识） ===
\setlength{\parindent}{0pt}
\setlength{\parskip}{0.35em plus 0.10em minus 0.08em}
\setlength{\emergencystretch}{3em}
\setlist{itemsep=0.25em, topsep=0.40em, parsep=0pt}
\allowdisplaybreaks[2]
% === 错题本排版优化结束 ===
\begin{document}
\section*{上海交通大学2025年高等代数真题}
\section*{36．上海交通大学2025年高等代数试卷（回忆版）}

\subsection{题目1}
1．证明：\(x^{p}+p x-p+1\) 在有理数域上不可约，其中 \(p\) 为奇素数。

\subsection{题目2}
2．求矩阵

\[
A=\left(\begin{array}{cccc}
2 & 0 & 3 & 7 \\
2 & 0 & -3 & -7 \\
2 & 0 & -3 & 7 \\
2 & 0 & -3 & -7
\end{array}\right)
\]

的行空间，列空间，零空间，左零空间的一组基（此处原本为三个矩阵相乘）。

\subsection{题目3}
3．设 \(A=\left(\begin{array}{lll}4 & 1 & 1 \\ 1 & 4 & 1 \\ 1 & 1 & 4\end{array}\right)\) ．

\begin{itemize}
  \item[（1）] 求 \(A\) 的特征多项式，特征值．
  \item[（2）] 求矩阵 \(P\) ，使得 \(P^{\mathrm{T}} A P\) 为对角阵。
  \item[（3）] 将 \(4 x_{1}^{2}+4 x_{2}^{2}+4 x_{3}^{2}+2 x_{1} x_{2}+2 x_{1} x_{3}+2 x_{2} x_{3}\) 化为标准形．
\end{itemize}

\subsection{题目4}
4．设 \(N=\left(\begin{array}{rrr}0 & 1 & \\ & 0 & \\ & & 1 \\ & & 0\end{array}\right)_{10 \times 10}\)

\begin{itemize}
  \item[（1）] 求 \(N^{3}\) 的若尔当标准形 \(J\) 。
  \item[（2）] 求置换矩阵 \(P\) 使得 \(P^{-1} N^{3} P=J, P\) 用标准单位列向量 \(e_{1}, e_{2}, \cdots, \varepsilon_{10}\) 表示。
  \item[（3）] 对于一般的 \(n\) 阶若尔当块 \(J_{n}(0)\) 和任意的正整数 \(k\) ，考虑 \(J_{n}^{k}(0)\) 的若尔当标准形．
\end{itemize}

\subsection{题目5}
5．证明：\(A^{2}=I\) 的充要条件是 \(r(A+I)+r(A-I)=n\) 。

\subsection{题目6}
6．设 \(\varphi\) 是 \(n\) 维西空间 \(V\) 上的线性变换，证明：\(\varphi\) 是正规算子的充要条件是存在复系数多项式 \(f(x)\) ，使得

\[
\varphi^{\bullet}=f(\varphi) .
\]

\subsection{题目7}
7．在 \(n\) 维欧氏空间中，两两内积为负的向量的个数至多是多少？请证明．

\subsection{题目8}
8．设 \(V, W\) 为线性空间（有限或无限维），\(V^{\prime}, W^{\prime}\) 为 \(V\) 和 \(W\) 的对偶空间，\(T\) 与 \(T^{\prime}\) 为线性算子，且\(T \in L(V, W), T^{\prime} \in L\left(W^{\prime}, V^{\prime}\right)\) 。证明：\(T\) 为满射的充要条件是 \(T^{\prime}\) 为单射。

\subsection{题目9}
9．设 \(f(x)=x^{n}+a_{1} x^{n-1}+\cdots+a_{n}, \mu\) 为 \(f(x)\) 的根，证明：\(|\mu| \leq \max \left\{1, \sum_{j=1}^{n}\left|a_{j}\right|\right\}\) ，且

\[
|\mu| \leq \max \left\{\left\{1+\left|a_{j}\right|, j=1,2, \cdots, n\right\} \cup\left\{\left|a_{n}\right|\right\}\right\} .
\]
\end{document}
